\documentclass{report}
\usepackage{luatexja} % LuaTeXで日本語を使うためのパッケージ
\usepackage{luatexja-fontspec} % LuaTeX用の日本語フォント設定

% --- 数学関連 ---
\usepackage{amsmath, amssymb, amsfonts, mathtools, bm, amsthm} % 基本的な数学パッケージ
\usepackage{type1cm, upgreek} % 数式フォントとギリシャ文字
\usepackage{physics, mhchem} % 物理や化学の記号や式の表記を簡単にする

% --- 表関連 ---
\usepackage{multirow, longtable, tabularx, array, colortbl, dcolumn, diagbox} % 表のレイアウトを柔軟にする
\usepackage{tablefootnote, truthtable} % 表中に注釈を追加、真理値表
\usepackage{tabularray} % 高度な表組みレイアウト

% --- グラフィック関連 ---
\usepackage{tikz, graphicx} % 図の描画と画像の挿入
\usepackage{background} % ウォーターマークの設定
\usepackage{caption, subcaption} % 図や表のキャプション設定
\usepackage{float, here} % 図や表の位置指定

% --- レイアウトとページ設定 ---
\usepackage{fancyhdr} % ページヘッダー、フッター、余白の設定
\usepackage[top = 20truemm, bottom = 20truemm, left = 20truemm, right = 20truemm]{geometry}
\usepackage{fancybox, ascmac} % ボックスのデザイン

% --- 色とスタイル ---
\usepackage{xcolor, color, colortbl, tcolorbox} % 色とカラーボックス
\usepackage{listings, jvlisting} % コードの色付けとフォーマット

% --- 参考文献関連 ---
\usepackage{biblatex, usebib} % 参考文献の管理と挿入
\usepackage{url, hyperref} % URLとリンクの設定

% --- その他の便利なパッケージ ---
\usepackage{footmisc} % 脚注のカスタマイズ
\usepackage{multicol} % 複数段組
\usepackage{comment} % コメントアウトの拡張
\usepackage{siunitx} % 単位の表記
\usepackage{docmute}
% \usepackage{appendix}
% --- tcolorboxとtikzの設定 ---
\tcbuselibrary{theorems, breakable} % 定理のボックスと改ページ設定
\usetikzlibrary{decorations.markings, arrows.meta, calc} % tikzの装飾や矢印の設定

% --- 定理スタイルと数式設定 ---
\theoremstyle{definition} % 定義スタイル
\numberwithin{equation}{section} % 式番号をサブセクション単位でリセット

% --- hyperrefの設定 ---
\hypersetup{
  setpagesize = false,
  bookmarks = true,
  bookmarksdepth = tocdepth,
  bookmarksnumbered = true,
  colorlinks = false,
  pdftitle = {}, % PDFタイトル
  pdfsubject = {}, % PDFサブジェクト
  pdfauthor = {}, % PDF作者
  pdfkeywords = {} % PDFキーワード
}

% --- siunitxの設定 ---
\sisetup{
  table-format = 1.5, % 小数点以下の桁数
  table-number-alignment = center, % 数値の中央揃え
}

% --- 透かし画像の設定 --- 
\backgroundsetup{
  scale=0.5,                       % 画像のスケール
  % color=black,                   % 画像の色（透かし用に半透明が推奨）
  opacity=0.2,                   % 透かしの透明度（0が完全透明、1が完全不透明）
  angle=0,                       % 画像の角度
  position = current page.south east,  % ページの右下
  hshift=-6cm, % 右方向へのシフト（負の値で内側に移動）
  vshift=5cm,  % 上方向へのシフト（正の値で内側に移動）
  contents={\includegraphics{./fig/appilogo-circular-full.png}} % 画像のパス
}

% --- その他の設定 ---
\allowdisplaybreaks % 数式の途中改ページ許可
\newcolumntype{t}{!{\vrule width 0.1pt}} % 新しいカラムタイプ
\newcolumntype{b}{!{\vrule width 1.5pt}} % 太いカラム
\UseTblrLibrary{amsmath, booktabs, counter, diagbox, functional, hook, html, nameref, siunitx, varwidth, zref} % tabularrayのライブラリ
\setlength{\columnseprule}{0.4pt} % カラム区切り線の太さ
\captionsetup[figure]{font = bf} % 図のキャプションの太字設定
\captionsetup[table]{font = bf} % 表のキャプションの太字設定
\captionsetup[lstlisting]{font = bf} % コードのキャプションの太字設定
\captionsetup[subfigure]{font = bf, labelformat = simple} % サブ図のキャプション設定
\setcounter{secnumdepth}{4} % セクションの深さ設定
\newcolumntype{d}{D{.}{.}{5}} % 数値のカラム
\newcolumntype{M}[1]{>{\centering\arraybackslash}m{#1}} % センター揃えのカラム
\DeclareMathOperator{\diag}{diag}
\everymath{\displaystyle} % 数式のスタイル
\newcommand{\inner}[2]{\left\langle #1, #2 \right\rangle}
\renewcommand{\figurename}{図}
\renewcommand{\i}{\mathrm{i}} % 複素数単位i
\renewcommand{\laplacian}{\grad^2} % ラプラシアンの記号
\renewcommand{\thesubfigure}{(\alph{subfigure})} % サブ図の番号形式
\newcommand{\m}[3]{\multicolumn{#1}{#2}{#3}} % マルチカラムのショートカット
\renewcommand{\r}[1]{\mathrm{#1}} % mathrmのショートカット
\newcommand{\e}{\mathrm{e}} % 自然対数の底e
\newcommand{\Ef}{E_{\mathrm{F}}} % フェルミエネルギー
\renewcommand{\c}{\si{\degreeCelsius}} % 摂氏記号
\renewcommand{\d}{\r{d}} % d記号
\renewcommand{\t}[1]{\texttt{#1}} % タイプライタフォント
\newcommand{\kb}{k_{\mathrm{B}}} % ボルツマン定数
% \renewcommand{\phi}{\varphi} % ϕをφに変更
\renewcommand{\epsilon}{\varepsilon}
\newcommand{\fullref}[1]{\textbf{\ref{#1} \nameref{#1}}}
\newcommand{\reff}[1]{\textbf{図\ref{#1}}} % 図参照のショートカット
\newcommand{\reft}[1]{\textbf{表\ref{#1}}} % 表参照のショートカット
\newcommand{\refe}[1]{\textbf{式\eqref{#1}}} % 式参照のショートカット
\newcommand{\refp}[1]{\textbf{コード\ref{#1}}} % コード参照のショートカット
\renewcommand{\lstlistingname}{コード} % コードリストの名前
\renewcommand{\theequation}{\thesection.\arabic{equation}} % 式番号の形式
\renewcommand{\footrulewidth}{0.4pt} % フッターの線
\newcommand{\mar}[1]{\textcircled{\scriptsize #1}} % 丸囲み文字
\newcommand{\combination}[2]{{}_{#1} \mathrm{C}_{#2}} % 組み合わせ
\newcommand{\thline}{\noalign{\hrule height 0.1pt}} % 細い横線
\newcommand{\bhline}{\noalign{\hrule height 1.5pt}} % 太い横線

% --- カスタム色定義 ---
\definecolor{burgundy}{rgb}{0.5, 0.0, 0.13} % バーガンディ色
\definecolor{charcoal}{rgb}{0.21, 0.27, 0.31} % チャコール色
\definecolor{forest}{rgb}{0.0, 0.35, 0} % 森の緑色

% --- カスタム定理環境の定義 ---
\newtcbtheorem[number within = chapter]{myexc}{練習問題}{
  fonttitle = \gtfamily\sffamily\bfseries\upshape,
  colframe = forest,
  colback = forest!2!white,
  rightrule = 1pt,
  leftrule = 1pt,
  bottomrule = 2pt,
  colbacktitle = forest,
  theorem style = standard,
  breakable,
  arc = 0pt,
}{exc-ref}
\newtcbtheorem[number within = chapter]{myprop}{命題}{
  fonttitle = \gtfamily\sffamily\bfseries\upshape,
  colframe = blue!50!black,
  colback = blue!50!black!2!white,
  rightrule = 1pt,
  leftrule = 1pt,
  bottomrule = 2pt,
  colbacktitle = blue!50!black,
  theorem style = standard,
  breakable,
  arc = 0pt
}{proposition-ref}
\newtcbtheorem[number within = chapter]{myrem}{注意}{
  fonttitle = \gtfamily\sffamily\bfseries\upshape,
  colframe = yellow!20!black,
  colback = yellow!50,
  rightrule = 1pt,
  leftrule = 1pt,
  bottomrule = 2pt,
  colbacktitle = yellow!20!black,
  theorem style = standard,
  breakable,
  arc = 0pt
}{remark-ref}
\newtcbtheorem[number within = chapter]{myex}{例題}{
  fonttitle = \gtfamily\sffamily\bfseries\upshape,
  colframe = black,
  colback = white,
  rightrule = 1pt,
  leftrule = 1pt,
  bottomrule = 2pt,
  colbacktitle = black,
  theorem style = standard,
  breakable,
  arc = 0pt
}{example-ref}
\newtcbtheorem[number within = chapter]{exc}{Requirement}{myexc}{exc-ref}
\newcommand{\rqref}[1]{{\bfseries\sffamily 練習問題 \ref{exc-ref:#1}}}
\newtcbtheorem[number within = chapter]{definition}{Definition}{mydef}{definition-ref}
\newcommand{\dfref}[1]{{\bfseries\sffamily 定義 \ref{definition-ref:#1}}}
\newtcbtheorem[number within = chapter]{prop}{命題}{myprop}{proposition-ref}
\newcommand{\prref}[1]{{\bfseries\sffamily 命題 \ref{proposition-ref:#1}}}
\newtcbtheorem[number within = chapter]{rem}{注意}{myrem}{remark-ref}
\newcommand{\rmref}[1]{{\bfseries\sffamily 注意 \ref{remark-ref:#1}}}
\newtcbtheorem[number within = chapter]{ex}{例題}{myex}{example-ref}
\newcommand{\exref}[1]{{\bfseries\sffamily 例題 \ref{example-ref:#1}}}
% --- 再定義コマンド ---
% \mathtoolsset{showonlyrefs=true} % 必要な式番号のみ表示
\pagestyle{fancy} % ヘッダー・フッターのスタイル設定
\chead{応用量子物性講義ノート} % 中央ヘッダー
% \rhead{}
\fancyhead[R]{\rightmark}
\renewcommand{\sectionmark}[1]{\markright{\thesection\ #1}}
\cfoot{\thepage} % 中央フッターにページ番号
\lhead{}
\rfoot{Yuto Masuda and Haruki Aoki} % 右フッターに名前
\setcounter{tocdepth}{4} % 目次の深さ
\makeatletter
\@addtoreset{equation}{section} % サブセクションごとに式番号をリセット
\makeatother

% --- メタ情報 ---
\title{応用量子物性講義ノート}
\date{更新日\today}
\author{Yuto Masuda and Haruki Aoki}

\begin{document}
  \textbf{スピン軌道相互作用}(spin-orbit coupling)は\textbf{スピントロニクス}(spintronics)において非常に重要な役割を担っている．
  スピン軌道相互作用によるハミルトニアンは
  \begin{align}
    \hat{H}_{\r{SO}} = \lambda_{\r{SO}} \hat{\bm{L}} \cdot \hat{\bm{S}}
  \end{align}
  で表される．

  これは定性的には次のように導出される．原子核の周りを動く電子を考える．静止系から見ると電子は電場$\bm{E} = -\frac{\partial V(r)}{\partial t}\frac{\bm{r}}{r}$を感じている．
  これを電子と共に動く慣性系で見る．この系では電子は静止しており，その周りを原子核が回っている．荷電粒子の運動は磁場を発生させるため，電子は磁場$\bm{B}$を感じる．この磁場はLorentz変換により
  \begin{align}
    \bm{B} = -\frac{\bm{v}}{c} \times \bm{E}
  \end{align}
  で電場と結ばれている．また，スピンと磁場は以下のように相互作用する．
  \begin{align}
    -\bm{\mu}\cdot\bm{B} \propto -\bm{S}\cdot\bm{B}
  \end{align}
  以上よりスピン軌道相互作用が得られる．
  \begin{align}
    H_{\r{SO}} = -\bm{\mu}\cdot\bm{B} &\propto -\bm{S}\cdot\bm{B}\\
    &= \bm{S} \cdot (\frac{\bm{v}}{c} \times \bm{E})\\
    &= \bm{S} \cdot(\frac{\bm{v}}{c} \times -\frac{\partial V(r)}{\partial t}\frac{\bm{r}}{r})\\
    &= \lambda_{\r{SO}} \bm{S}\cdot(\bm{r} \times \bm{p})\\
    &= \lambda_{\r{SO}} \bm{S} \cdot \bm{L}
  \end{align}
  つまり，電子の感じる仮想的な磁場がスピン軌道相互作用の起源である．

  Dirac方程式を用いてスピン軌道相互作用を導出する．Dirac方程式は以下である．
  \begin{align}
    \i\hbar\frac{\partial}{\partial t}\psi = (c\bm{\alpha} \cdot \bm{p} + \beta mc^2 + e\phi) \psi
  \end{align}
  前節と同様に
  \begin{align}
    \psi=
    \begin{pmatrix}
      \psi_a\\
      \psi_b
    \end{pmatrix}
  \end{align}
  とする．さらにエネルギー演算子を
  \begin{align}
    \hat{E} \equiv \i\hbar\frac{\partial}{\partial t} \equiv \hat{E}' + mc^2
  \end{align}
  と置く．最右辺ではエネルギーを静止エネルギーと非相対論的エネルギーに分けた．電磁場を$\bm{A} =0$として加え，前節と同様に計算を進める．
  \begin{align}
    (\hat{E}' + mc^2) 
    \begin{pmatrix}
      \psi_a\\
      \psi_b
    \end{pmatrix}
    = c\bm{\sigma}\cdot \hat{\bm{p}}
    \begin{pmatrix}
      \psi_b\\
      \psi_a
    \end{pmatrix}
    mc^2
    \begin{pmatrix}
      \psi_a\\
      -\psi_b
    \end{pmatrix}
    + e\phi
    \begin{pmatrix}
      \psi_a\\
      \psi_b
    \end{pmatrix}
  \end{align}
  上式より連立方程式
  \begin{align}
    \begin{dcases}
      (\hat{E}' - e\phi) \psi_a - c\bm{\sigma}\cdot\hat{\bm{p}} \psi_b=0\\
      (\hat{E}' + 2mc^2 - e\phi) \psi_b - c \bm{\sigma}\cdot\bm{p}\psi_a = 0
    \end{dcases}
  \end{align}
  を得る．2式目を1式目に代入する．
  \begin{align}
    (\hat{E}' - e\phi) \psi_a &= (c\bm{\sigma}\cdot\hat{\bm{p}})(\hat{E}' + 2mc^2 - e\phi)^{-1}c\bm{\sigma}\cdot\hat{\bm{p}}\psi_a\\
    &\simeq \frac{1}{2m}(\bm{\sigma\cdot\hat{\bm{p}}})\qty(1 - \frac{\hat{E}' - e\phi}{2mc^2})(\bm{\sigma}\cdot\hat{\bm{p}})\psi_a\\
  \end{align}
  ここで，
  \begin{align}
    \hat{\bm{p}}(e\phi) = (e\phi)\hat{\bm{p}} - \i e \hbar \grad\phi
  \end{align}
  であるため，さらに変形を進めて，
  \begin{align}
    (\hat{E}' - e\phi) \psi_a &= \frac{1}{2m}\bm{\sigma}\cdot\qty(\hat{\bm{p}} - \frac{\hat{E}' - (e\phi)\hat{\bm{p}} + \i e\hbar \grad\phi}{2mc^2})(\bm{\sigma}\cdot\hat{\bm{p}}) \psi_a\\
    &= \qty[\frac{1}{2m}(\bm{\sigma}\cdot\hat{\bm{p}})(\bm{\sigma}\cdot\hat{\bm{p}}) - \frac{\hat{E}' - e\phi}{4m^2c^2}(\bm{\sigma}\cdot\hat{\bm{p}})(\bm{\sigma}\cdot\hat{\bm{p}}) - \i \frac{e\hbar}{4m^2c^2}(\bm{\sigma}\cdot\grad\phi)(\bm{\sigma}\cdot\hat{\bm{p}})]\psi_a\\
    &= \qty[\qty(1 - \frac{\hat{E}' - e\phi}{2mc^2})\frac{\hat{\bm{p}}^2}{2m} - \i \frac{e\hbar}{4m^2c^2}(\grad\phi)\cdot\hat{\bm{p}} + \frac{e\hbar}{4m^2c^2}\bm{\sigma}\cdot(\grad\phi\times\hat{\bm{p}})]\psi_a
  \end{align}
  が得られる．4項目
  \begin{align}
    \frac{e\hbar}{4m^2c^2}\bm{\sigma}\cdot(\grad\phi\times\hat{\bm{p}})
  \end{align}
  がスピン軌道相互作用を表している．特に，球対称ポテンシャルのとき，
  \begin{align}
    \frac{e\hbar}{4m^2c^2}\bm{\sigma}\cdot(\grad\phi\times\hat{\bm{p}}) &= \frac{e\hbar}{4m^2c^2}\frac{2}{\hbar}\bm{\bm{S}}\cdot(\frac{1}{r}\dv{\phi}{r}\bm{r}\times\hat{\bm{p}})\\
    &= \frac{e}{2m^2c^2r}\dv{\phi}{r}\hat{\bm{S}}\cdot\hat{\bm{L}}
  \end{align}
  となり，スピン角運動量と軌道角運動量がカップリングしていることがわかる．

  次に，スピン軌道相互作用の大きさを考える．全角運動量を$\bm{J}$として
  \begin{align}
    \bm{J}^2 = (\bm{L} + \bm{S})^2 = \bm{L}^2 + \bm{S}^2 + 2\bm{L}\cdot \bm{S}\\
    \bm{L}\cdot\bm{S} = \bm{J^2} - \bm{L}^2 - \bm{S}^2
  \end{align}
  だから$\hat{\bm{L}}\cdot\hat{\bm{S}}$の期待値は
  \begin{align}
    \langle\hat{\bm{L}}\cdot\hat{\bm{S}} \rangle = \frac{1}{2}\qty[j(j+1) - l(l+1) - \frac{1}{2}\qty(\frac{1}{2}+1)]\hbar^2
  \end{align}
  である．$j = l \pm 1/2$だから，
  $\langle\hat{\bm{L}}\cdot\hat{\bm{S}} \rangle = \frac{1}{2}l\hbar^2,-\frac{1}{2}(l+1)\hbar^2$である．これを$\nu\hbar^2$とおく．
  $\phi(r) = -Ze/r$($Z$は原子番号)とすると，
  \begin{align}
    \langle \frac{e}{2m^2c^2r}\dv{\phi}{r}\hat{\bm{S}}\cdot\hat{\bm{L}} \rangle &= \nu\hbar^2 \langle \frac{e}{2m^2c^2r}\dv{\phi}{r} \rangle\\
    &= \mu\hbar^2 \int_{0}^{\infty} r^2 \dd{r} \abs{R_{nl}}^2 \qty(\frac{e}{2m^2c^2r}\dv{\phi}{r})\\
    &= \frac{\mu\hbar^2 Z e^2}{2m^2c^2}\int_{0}^{\infty} \frac{1}{r}R_{nl}^2 \dd{r}\\
    &= \frac{\mu\hbar^2 Z^4 e^2}{2m^2 c^2 a_0^3n^3l(l+1/2)(l+1)}
  \end{align}
  がスピン軌道相互作用の期待値である．ここで，$a_0=\hbar^2/(\mu e^2),\mu = Mm_e(M+m_e)$である．重要なのはスピン軌道相互作用の大きさは
  原子番号$Z$の4乗に比例することである．つまり，周期表の下のにある重い元素ほどスピン軌道相互作用が強いことがわかる．
\end{document}